\documentclass[12pt]{article}
\usepackage[utf8]{inputenc}
\usepackage[T1]{fontenc}
\usepackage{hyperref}
\title{Waru Language Resources for K-12 Curriculum Development}
\author{Compiled by Claude Sonnet 5.5 (claude-sonnet-5-5) for Andr\'as Kornai}
\date{}
\begin{document}
\maketitle

\section{Introduction}
Waru (also Mopute or Mapute; ISO 639-3 wru; Glottocode waru1266) is an endangered Austronesian language of Sulawesi, Indonesia, in the Bungku-Tolaki group. Sources place it in Southeast Sulawesi (Mopute village, North Konawe). Joshua Project gives a conflicting Central Sulawesi location, which I treat as doubtful. Estimates are about 350 speakers. Little exists for it online.

\section{Primary Educational Resources}
No primers, textbooks or school materials were found in search results.

\section{Linguistic and Reference Materials}
\begin{itemize}
\item Wikipedia entry for the Waru language: \url{https://en.wikipedia.org/wiki/Waru_language}
\item Word-list data in the ASJP database: \url{https://asjp.clld.org/languages/WARU}
\item Endangered Languages Project profile: \url{https://endangeredlanguages.com/lang/2802}
\end{itemize}
No dictionary or grammar was found. Tukang Besi, which also came up in the search, is a different language and is not a Waru resource.

\section{Religious and Community Resources}
Joshua Project has a people-group page for the Waru in Indonesia, with a population of 600: \url{https://joshuaproject.net/people_groups/15875}. No scripture or community organisation was found.

\section{Digital and Audio Resources}
Verbix has a short page: \url{https://docs.verbix.com/Languages/Waru}. No audio, video or software was found.

\section{Curriculum Development Context}
Almost nothing exists to build on beyond wordlists and general descriptions. A curriculum would need fresh documentation with the community, and the actual location of the speakers should be confirmed first.
\end{document}